\begin{example}[Chain Rule learning in FOL]

    Let us consider a set of rules, which are formulas
    \begin{align*}
        q_0\land \ldots\land q_{n-1} \implies q_n
    \end{align*}
    where each $q_i$ is a FOL formula depending on variables $\indvariableof{i}$ and $\indvariableof{i+1}$

    We select each $q_i$ from a set $\arbsetof{i}$ enumerated by $\indvariableof{i}$.


    Effective parametrization: Sometimes we can contract the selector cores with the candidate atom cores to increase efficiency.

    \stantikzpic{
        \begin{scope}[shift={(-2,-2)}]
            \drawtdwire{-7}{3}{$\headvariable$}
            \drawhblock{-7}{0}{2}{$\bencodingof{\exformula}$}
            \drawldwire{-5}{0}{$\selvariable$}
            \drawtdwire{-8.5}{-1}{$\indvariableof{0}$}
            \drawtdwire{-5.5}{-1}{$\indvariableof{1}$}

            \drawtextnode{-1}{0}{$=$}
        \end{scope}


        \drawtdwire{5}{5}{$\headvariable$}
        \drawhblock{5}{2}{2}{$\bencodingof{\vselectionmap}$}
        \drawldwire{7}{2}{$\selvariable$}

        \draw [->-] (1,-1) to [bend right = -10] (3.5,1);
        \drawhblock{1}{-2}{2}{$\bencodingof{\folpredicateof{0}}$}

        \draw [->-] (9,-1) to [bend right = 10] (6.5,1);
        \drawhblock{9}{-2}{2}{$\bencodingof{\folpredicateof{\seldim-1}}$}

        \draw[->-] (3,-5) to [bend right = -15] (-0.5,-3);
        \draw[->-] (3,-5) to [bend right = 15] (7.5,-3);
        \drawvariabledot{3}{-5}
        \drawtdwire{3}{-5}{$\indvariableof{0}$}

        \draw[->-] (7,-5) to [bend right = -15] (2.5,-3);
        \draw[->-] (7,-5) to [bend right = 15] (10.5,-3);
        \drawvariabledot{7}{-5}
        \drawtdwire{7}{-5}{$\indvariableof{1}$}
    }


    Using the effective parameterization formulas by the tensor network:
    \stantikzpic{

        \drawtdwire{0}{3}{}
        \drawhblock{0}{0}{2}{$\bencodingof{\formulaof{0}}$}
        \draw[-<-] (2,0) to [bend right = -30] (3,-2);
        \drawtextnode{2.75}{0.25}{\colorlabelsize $\selvariableof{0}$}

        \draw[->-] (-1.5,-4) -- (-1.5,-1);
        \drawtdwire{-1.5}{-4}{}
        \drawtextnode{-2.25}{-5}{\colorlabelsize $\indvariableof{0}$}
        \drawvariabledot{-1.5}{-4}
        \draw[->-] (-1.5,-4) to [bend right=70] (21.5,-1);
%        \drawtdwire{3}{-4}{$\indvariableof{1}$}

        \draw[->-] (3,-4) to [bend right = -20] (1.5,-1);
        \drawvariabledot{3}{-4}
        \drawtdwire{3}{-4}{$\indvariableof{1}$}
        \draw[->-] (3,-4) to [bend right = 20] (4.5,-1);


        \drawtdwire{6}{3}{}
        \drawhblock{6}{0}{2}{$\bencodingof{\formulaof{1}}$}
        \drawtextnode{8.75}{0.25}{\colorlabelsize $\selvariableof{1}$}
        \draw[-<-] (8,0) to [bend right = -30] (9,-2);

        \draw[->-] (9,-4) to [bend right = -20] (7.5,-1);
        \drawvariabledot{9}{-4}
        \drawtdwire{9}{-4}{$\indvariableof{2}$}
%        \draw[->-] (9,-4) to [bend right = 20] (10.5,-1);

        \drawtextnode{10}{-1}{$\cdots$}

        \drawtextnode{10}{-3}{$\cdots$}

        %\draw[->-] (11,-4) to [bend right = -20] (9.5,-1);
        \drawvariabledot{11}{-4}
        \drawtdwire{11}{-4}{$\indvariableof{n-1}$}
        \draw[->-] (11,-4) to [bend right = 20] (12.5,-1);

        \draw[->-] (17,-4) to [bend right = -20] (15.5,-1);
        \drawvariabledot{17}{-4}
        \drawtdwire{17}{-4}{$\indvariableof{n}$}
        \draw[->-] (17,-4) to [bend right = 20] (18.5,-1);



        \drawtdwire{14}{3}{}
        \drawhblock{14}{0}{2}{$\bencodingof{\formulaof{n-1}}$}
        \drawtextnode{16.75}{0.25}{\colorlabelsize $\selvariableof{n-1}$}
        \draw[-<-] (16,0) to [bend right = -30] (17,-2);

        \drawhblock{20}{0}{2}{$\bencodingof{\formulaof{n}}$}
        \drawtextnode{22.75}{0.25}{\colorlabelsize $\selvariableof{n}$}
        \draw[-<-] (22,0) to [bend right = -30] (23,-2);

        \drawtdwire{7}{7}{}
        \drawhblock{7}{4}{8}{$\bencodingof{\land}$}

        \drawsmalltensor{14}{12}{$\tbasis$}
        \drawtdwire{14}{11}{}
        \drawhblock{14}{8}{8}{$\bencodingof{\impbincon}$}
        \drawtdwire{20}{7}{}
        \draw (20,1) -- (20,5);







    }



\end{example}